\documentclass[main.tex]{subfiles}
\begin{document}

Rising credit card use and rising merchant costs both follow from intense network competition channeled toward the wrong side of the market.
Because consumers are far more sensitive to rewards than merchants are to fees, networks compete for cardholders with generous rewards funded by merchant fees.
Price coherence ensures that these fees pass through to higher retail prices borne by all consumers, including those who pay with cash and debit. 
The rewards draw more consumers to credit cards, increasing merchants' costs.

The counterfactual results speak to contemporary legal and policy debates in payments and platform markets.
The Supreme Court's 2018 decision in \emph{Ohio v.\ American Express} held that plaintiffs must show net harm to the two-sided market as a whole, requiring evidence on both sides of the platform.
This paper provides a structural framework for carrying out that analysis, linking consumer rewards, merchant fees, and retail prices within a single equilibrium model of a two-sided card market.
Applied to credit card merchant fees, the framework shows that a modest cap reduces both fees and rewards yet raises total welfare by \$27 billion per year through more efficient payment choice.
% HARDCODED[table: welfare_table_baseline.tex, row: "Total", col: 1] current=27

The Durbin Amendment counterfactual shows that capping the wrong fees makes things worse.
Congress capped debit interchange in 2010 because it appeared supracompetitive, but those fees fund rewards that draw consumers away from high-fee credit cards.
Repealing the cap would raise welfare by \$7 billion per year,
% HARDCODED[table: welfare_table_baseline.tex, row: "Total", col: 3] current=7
making the current U.S.\ regime of capping debit but not credit fees worse than laissez-faire.

The merger counterfactual challenges the view that high merchant fees reflect too little competition.
In fact, it's precisely network competition that encourages the adoption of high-merchant-fee payment methods.
A monopoly network raises per-transaction fees, which is the harm that antitrust cases have targeted.
Yet aggregate merchant costs fall because the monopolist has no rival to outbid for cardholders, so it cuts rewards and credit use declines.
The merger increases total welfare by \$15 billion
% HARDCODED[table: welfare_table_baseline.tex, row: "Total", col: 6] current=15
because eliminating rewards competition reduces the overuse of credit cards.
Buy Now, Pay Later is the latest technology to woo consumers with generous terms funded by high merchant fees \parencite{Berg.etal2024}.
Whether these services raise welfare depends on whether they shift competition toward merchants or simply open another front in the war for cardholders.

A broader lesson from the counterfactuals is that which side of the market networks compete for matters as much as how intensely they compete.
The dual routing counterfactual illustrates this. 
When consumer multi-homing forces networks to compete for merchants rather than cardholders, both credit card fees and rewards fall, and total welfare rises on net.
\end{document}
